\documentclass[]{beamer}
%\documentclass[handout]{beamer}
% ***************************************************************
% for handout, change only this...
%   \documentclass[twocolumn]{article}
%   \usepackage{beamerarticle}
%   \setlength{\textwidth}{7.5in}
%   \setlength{\textheight}{9.8in}
%   \setlength{\topmargin}{-1in}
%   \setlength{\oddsidemargin}{-.52in}
%   \setlength{\evensidemargin}{-.52in}

%\usepackage{beamerprosper}
%\usetheme{Warsaw}
%\usecolortheme{orchid}

\usepackage{graphicx}
\usepackage{amsmath,amssymb,array,comment,eucal}
\usepackage{xcolor}
\definecolor{beamer@blendedblue}{RGB}{86,155,189}
\definecolor{myblue}{RGB}{12,76,138}
\setbeamercolor{structure}{fg=myblue}
\definecolor{Ftitle}{RGB}{12,76,138}
\definecolor{Descitem}{RGB}{238,238,244}
\definecolor{StdTitle}{RGB}{12,76,138}
\definecolor{StdBody}{RGB}{213,24,0}
\definecolor{StdBody}{RGB}{213,24,0}

\definecolor{AlTitle}{RGB}{255, 190, 190}
\definecolor{AlBody}{RGB}{213,24,0}

\definecolor{ExTitle}{RGB}{201, 217, 217}
\definecolor{ExBody}{RGB}{213,24,0}

\setbeamercolor{frametitle}{fg = Ftitle}
\setbeamercolor{title}{fg = Ftitle}
\setbeamercolor{item}{fg = Ftitle}
\setbeamercolor{subitem}{fg = Ftitle}
\setbeamercolor{subsubitem}{fg = Ftitle}
\setbeamercolor{description item}{fg = myblue}
\setbeamercolor{titlelike}{fg=myblue}

\DeclareMathOperator{\sgn}{sgn}
\newcommand{\e}{\mathbf{e}}
\renewcommand{\P}{\mathbf{P}}
\newcommand{\F}{\mathbf{F}}
\newcommand{\R}{\textsf{R}}
\newcommand{\mat}[1] {\mathbf{#1}}
%\newcommand{\ind}{\mathrel{\mathop{\sim}\limits^{\mathit{ind}}}}
%\newcommand{\iid}{\mathrel{\mathop{\sim}\limits^{\mathit{iid}}}}
\newcommand{\E}{\textsf{E}}
\newcommand{\SE}{\textsf{SE}}
\newcommand{\SSE}{\textsf{SSE}}
\newcommand{\RSS}{\textsf{RSS}}
\newcommand{\FSS}{\textsf{FSS}}
\renewcommand{\SS}{\textsf{SS}}
\newcommand{\MSE}{\textsf{MSE}}
\newcommand{\SSR}{\textsf{SSR}}
\newcommand{\Be}{\textsf{Beta}}
\newcommand{\St}{\textsf{St}}
\newcommand{\Ca}{\textsf{C}}
\newcommand{\Exp}{\textsf{Exp}}
\newcommand{\GDP}{\textsf{GDP}}
\newcommand{\NcSt}{\textsf{NcSt}}
\newcommand{\Bin}{\textsf{Bin}}
\newcommand{\NB}{\textsf{NegBin}}
\renewcommand{\NG}{\textsf{NG}}
\newcommand{\N}{\textsf{N}}
\newcommand{\Ber}{\textsf{Ber}}
\newcommand{\Poi}{\text{Poi}}
\newcommand{\Gam}{\textsf{Gamma}}
\newcommand{\BB}{\textsf{BB}}
\newcommand{\Gm}{\textsf{G}}
\newcommand{\Un}{\textsf{Unif}}
\newcommand{\Ex}{\textsf{Exp}}
\newcommand{\DE}{\textsf{DE}}
\newcommand{\tr}{\textsf{tr}}
\newcommand{\cF}{{\cal{F}}}
\newcommand{\cL}{{\cal{L}}}
\newcommand{\cI}{{\cal{I}}}
\newcommand{\cB}{{\cal{B}}}
\newcommand{\cP}{{\cal{P}}}
\newcommand{\bbR}{\mathbb{R}}
\newcommand{\bbN}{\mathbb{N}}
\newcommand{\pperp}{\mathrel{{\rlap{$\,\perp$}\perp\,\,}}}
\newcommand{\OFP}{(\Omega,\cF, \P)}
\newcommand{\eps}{\boldsymbol{\epsilon}}
\newcommand{\1}{\mathbf{1}_n}
\newcommand{\gap}{\vspace{8mm}}
\newcommand{\ind}{\mathrel{\mathop{\sim}\limits^{\rm ind}}}
\newcommand{\simiid}{\ensuremath{\mathrel{\mathop{\sim}\limits^{\rm
iid}}}}
\newcommand{\eqindis}{\ensuremath{\mathrel{\mathop{=}\limits^{\rm D}}}}
\newcommand{\iid}{\textit{i.i.d.}}
\newcommand{\SSZ}{S_{zz}}
\newcommand{\SZW}{S_{zw}}
\newcommand{\Var}{\textsf{Var}}
\newcommand{\corr}{\textsf{corr}}
\newcommand{\diag}{\textsf{diag}}
\newcommand{\var}{\textsf{var}}
\newcommand{\Cov}{\textsf{Cov}}
\newcommand{\Sam}{{\cal S}}
\def\H{\mathbf{H}}
\newcommand{\I}{\mathbf{I}}
\newcommand{\Y}{\mathbf{Y}}
\newcommand{\tY}{\tilde{\mathbf{Y}}}
\newcommand{\Yhat}{\hat{\mathbf{Y}}}
\newcommand{\Yobs}{\mathbf{Y}_{{\cal S}}}
\newcommand{\barYobs}{\bar{Y}_{{\cal S}}}
\newcommand{\barYmiss}{\bar{Y}_{{\cal S}^c}}
\def\bv{\mathbf{b}}
\def\X{\mathbf{X}}
\def\tX{\tilde{\mathbf{X}}}
\def\x{\mathbf{x}}
\def\xbar{\bar{\mathbf{x}}}
\def\Xbar{\bar{\mathbf{X}}}
\def\Xg{\mathbf{X}_{\boldsymbol{\gamma}}}
\def\Ybar{\bar{\Y}}
\def\ybar{\bar{y}}
\def\y{\mathbf{y}}
\def\Yf{\mathbf{Y_f}}
\def\W{\mathbf{W}}
\def\L{\mathbf{L}}
\def\w{\mathbf{w}}
\def\U{\mathbf{U}}
\def\V{\mathbf{V}}
\def\Q{\mathbf{Q}}
\def\Z{\mathbf{Z}}
\def\z{\mathbf{z}}
\def\v{\mathbf{v}}
\def\u{\mathbf{u}}

\def\zero{\mathbf{0}}
\newcommand{\one}{\mathbf{1}}
\newcommand{\taub}{\boldsymbol{\tau}}
\newcommand{\betav}{\boldsymbol{\beta}}
\newcommand{\alphav}{\boldsymbol{\alpha}}
\newcommand{\A}{\mathbf{A}}
\def\a{\mathbf{a}}
\def\K{\mathbf{K}}
\newcommand{\B}{\mathbf{B}}
\def\b{\boldsymbol{\beta}}
\def\bhat{\hat{\boldsymbol{\beta}}}
\def\btilde{\tilde{\boldsymbol{\beta}}}
\def\tb{\boldsymbol{\theta}}
\def\bg{\boldsymbol{\beta_\gamma}}
\def\bgnot{\boldsymbol{\beta_{(-\gamma)}}}
\def\mub{\boldsymbol{\mu}}
\def\tmub{\tilde{\boldsymbol{\mu}}}
\def\muhat{\hat{\boldsymbol{\mu}}}
\def\tb{\boldsymbol{\theta}}
\def\tk{\boldsymbol{\theta}_k}
\def\tj{\boldsymbol{\theta}_j}
\def\Mk{\boldsymbol{{\cal M}}_k}
\def\M{\boldsymbol{{\cal M}}}
\def\Mj{\boldsymbol{{\cal M}}_j}
\def\Mi{\boldsymbol{{\cal M}}_i}
\def\Mg{{\boldsymbol{{\cal M}_\gamma}}}
\def\Mnull{\boldsymbol{{\cal M}}_{N}}
\def\gMPM{\boldsymbol{\gamma}_{\text{MPM}}}
\def\gHPM{\boldsymbol{\gamma}_{\text{HPM}}}
\def\Mfull{\boldsymbol{{\cal M}}_{F}}
\def\tg{\boldsymbol{\theta}_{\boldsymbol{\gamma}}}
\def\g{\boldsymbol{\gamma}}
\def\eg{\boldsymbol{\eta}_{\boldsymbol{\gamma}}}
\def\G{\mathbf{G}}
\def\cM{\cal M}
\def\D{\Delta}
\def \shat{{\hat{\sigma}}^2}
\def\uv{\mathbf{u}}
\def\l {\lambda}
\def\d{\delta}
\def\Sigmab{\boldsymbol{\Sigma}}
\def\Lambdab{\boldsymbol{\Lambda}}
\def\lambdab{\boldsymbol{\lambda}}
\def\Mg{{\cal M}_\gamma}
\def\S{{\cal{S}}}
\def\qg{p_{\boldsymbol{\gamma}}}
\def\pg{p_{\boldsymbol{\gamma}}}
%\def\t{\mathbf{t}}
\def\T{\boldsymbol{\Theta}}
\def\Tb{\boldsymbol{\Theta}}

\title{Bayesian Model Averaging}

\author{Hoff Chapter 9, Liang et al 2007, Hoeting et al (1999), Clyde \&
 George (2004) Statistical Science}
\date{\today}
%\Logo(-1.9,7.3){\includegraphics[width=.5in]{../eps/duke}}
% Optional: text to put in the bottom of each slide.
% By default, the title of the talk will be placed there.
%\slideCaption{\textit{October 28, 2005 }}

\usepackage{Sweave}
\begin{document}
\input{bma-concordance}
% make the title slide
\maketitle


\begin{frame}\frametitle{Prior Distributions}
  \begin{itemize}
  \item Bayesian Model choice requires  proper prior distributions on
    parameters that are not common across models  \pause
  \item Vague but proper priors may lead to paradoxes! \pause
  \item Conjugate Normal-Gammas lead to closed form expressions for
  marginal likelihoods, Zellner's g-prior is the most popular. \pause
  \end{itemize}
\end{frame}


\begin{frame}\frametitle{Prior Distributions}
  \begin{itemize}
  \item ``Spike and Slab'' - Lempers (1971) Mitchell \& Beauchamp (1988) \pause
  \item ``Spike and Bell''  Leamer (1978)  in BMA \pause
  \item   mixture of 2 normals - concentrated and dispersed - SSVS Gibbs Sampler -  George \& McCulloch (1993) \pause
  \item Back to ``Spike and Bell'' Hoeting, Raftery \& Madigan $MC^3$
    (1997) and George \& McCulloch (1997) collapsed MCMC after
    integrating out $\bg$
  \item Conjugate Normal-Gammas lead to closed form expressions for
  marginal likelihoods, Zellner's g-prior is the most popular. \pause
  \end{itemize}
\end{frame}

\begin{frame}\frametitle{Zellner's g-prior}
  Centered model:  $$\Y = \1 \alpha + \X^c \b + \epsilon$$
  where $\X^c$ is the centered design matrix where all variables have
  had their mean subtracted \pause
\begin{itemize}
\item   $p(\alpha) \propto 1$ \pause
\item  $p(\sigma^2) \propto 1/\sigma^2$  \pause
\item  $\b_\gamma \mid \alpha, \sigma^2, \g \sim \N(0, g \sigma^2
  ({\Xg^c}^\prime \Xg^c)^{-1})$ \pause
%\item  take $g=n$
\end{itemize}
which leads to marginal likelihood of $\Mg$ that is proportional
to $$ p(\Y \mid \Mg) = C (1 + g)^{\frac{n-p-1}{2}} ( 1 + g (1 -
 R^2_\gamma))^{- \frac{(n-1)}{2}}$$
where $R^2$ is the usual coefficient of determination for model $\Mg$.
\pause

Trade-off of model complexity versus goodness of fit

\bigskip
Lastly, assign distribution to space of models
\end{frame}

\begin{frame}{Sketch}
\begin{itemize}
  \item  Integrate out $\bg$  using sums of normals \pause
  \item  Find inverse of $\I_n + g \P_{\Xg}$  (properties of projections) \pause
  \item Find determinant of $\phi (\I_n + g \P_{\Xg})$  \pause
  \item Integrate out intercept (normal)  \pause
  \item Integrate out $\phi$  (gamma)  \pause
  \item algebra to simplify in from quadratic forms to  $R^2_{\g}$
\end{itemize}

\end{frame}
\begin{frame}
  \frametitle{Priors on Model Space}
  $p(\Mg) \Leftrightarrow p(\g)$
  \begin{itemize}
  \item $p(\gamma_j = 1) = .5 \Rightarrow P(\Mg) = .5^p$  Uniform on space of models \pause $\pg \sim \Bin(p, .5)$
\item $\gamma_j \mid \pi \simiid \Ber(\pi)$ and $\pi \sim \Be(a,b)$ then  $\pg \sim \BB_p(a, b)$
$$
p(\pg \mid p, a, b) = \frac{ \Gamma(p + 1) \Gamma(\pg + a) \Gamma(p - \pg + b) \Gamma (a + b) }{\Gamma(\pg+1) \Gamma(p - \pg + 1) \Gamma(p + a + b) \Gamma(a) \Gamma(b)}
$$
\item $\pg \sim \BB_p(1, 1) \sim \Un(0, p)$
  \end{itemize}
\end{frame}

\begin{frame}[fragile]
\frametitle{USair Data}
\begin{Schunk}
\begin{Sinput}
> library(BAS)
> data(usair, package="HH")
> poll.bma = bas.lm(log(SO2) ~ temp + log(mfgfirms) +
+                              log(popn) + wind +
+                              precip + raindays,
+                   data=usair,
+                   prior="g-prior",
+                   alpha=nrow(usair), # g = n
+                   n.models=2^6,
+                   modelprior = uniform(),
+                   method="deterministic")
\end{Sinput}
\end{Schunk}

% Marginal Posterior Inclusion Probabilities:
% Intercept  temp   log(mfgfirms)  log(popn)   wind  precip  raindays
%   1.0000   0.9755       0.7190     0.2757  0.7654  0.5994   0.3104

\end{frame}

\begin{frame}[fragile]\frametitle{Summary}

\begin{small}
\begin{Schunk}
\begin{Sinput}
> poll.bma
\end{Sinput}
\begin{Soutput}
Call:
bas.lm(formula = log(SO2) ~ temp + log(mfgfirms) + log(popn) + 
    wind + precip + raindays, data = usair, n.models = 2^6, prior = "g-prior", 
    alpha = nrow(usair), modelprior = uniform(), method = "deterministic")


 Marginal Posterior Inclusion Probabilities: 
    Intercept           temp  log(mfgfirms)      log(popn)           wind  
       1.0000         0.9755         0.7190         0.2757         0.7654  
       precip       raindays  
       0.5994         0.3104  
\end{Soutput}
\end{Schunk}
\end{small}

\end{frame}
\begin{frame}[fragile]\frametitle{Plots}

\centering
\begin{Schunk}
\begin{Sinput}
> par(mfrow=c(2,2))
> plot(poll.bma, ask=F)
\end{Sinput}
\end{Schunk}
\end{frame}

\begin{frame}\frametitle{Posterior Distribution  with Uniform Prior on Model Space}


\begin{Schunk}
\begin{Sinput}
> image(poll.bma, rotate=FALSE)
\end{Sinput}
\end{Schunk}


\end{frame}

\begin{frame}[fragile]\frametitle{Posterior Distribution  with BB(1,1) Prior on Model Space}


\begin{Schunk}
\begin{Sinput}
> poll.bb.bma = bas.lm(log(SO2) ~ temp + log(mfgfirms) +
+                                 log(popn) + wind +
+                                 precip + raindays,
+                      data=usair,
+                      prior="g-prior",
+                      alpha=nrow(usair),
+                      n.models=2^6,  #enumerate
+                      modelprior=beta.binomial(1,1))
\end{Sinput}
\end{Schunk}

\end{frame}

\begin{frame}\frametitle{BB(1,1) Prior on Model Space}


\begin{Schunk}
\begin{Sinput}
> image(poll.bb.bma, rotate=FALSE)
\end{Sinput}
\end{Schunk}

\end{frame}
\begin{frame}
\frametitle{Jeffreys Scale of Evidence}
\begin{itemize}
\item Bayes Factor = ratio of marginal likelihoods \pause

\item Posterior odds = Bayes Factor $\times$ Prior odds \pause

\end{itemize}


$B = BF[\M_0 : \Mg]$   and $1/B = BF[\Mg : \M_0]$ \pause

\vspace{14pt}

\begin{tabular}{|r|l|} \hline \hline
Bayes Factor & Interpretation \\ \hline
$B \geq 1$ & $H_0$ supported \\
$1 > B \geq 10^{-\frac{1}{2}} $ & minimal evidence against $H_0$ \\
$ 10^{- \frac{1}{2}} > B  \geq 10^{-1}$ & substantial evidence against $H_0$ \\
$ 10^{-1} > B  \geq 10^{-2}$ & strong evidence against $H_0$ \\
$ 10^{-2} > B $ & decisive evidence against $H_0$ \\ \hline \hline
\end{tabular}

\pause
\vspace{24pt}
in context of testing one hypothesis with equal prior odds
Kass \& Raftery (JASA 1996)
\end{frame}


\begin{frame}[fragile]
\frametitle{Coefficients}
\begin{Schunk}
\begin{Sinput}
>  beta = coef(poll.bma)
>  par(mfrow=c(2,3));  plot(beta, subset=2:7,ask=F)
\end{Sinput}
\end{Schunk}

\end{frame}

\begin{frame}
  \frametitle{Bartlett's Paradox}

The Bayes factor for comparing $\Mg$ to the null
model:
$$
 BF(\Mg : \M_0) =    (1 + g)^{(n - 1 - \pg)/2} (1 + g(1 - R^2))^{-(n-1)/2}
$$
\pause
\begin{itemize}
\item For fixed sample size $n$ and $R^2$, consider taking values of  $g$ that
  go to infinity  \pause
\item Increasing vagueness in prior \pause
\item What happens to BF as $g \to \infty$?

\end{itemize}





\end{frame}

\begin{frame}
  \frametitle{Information Paradox}

The Bayes factor for comparing $\Mg$ to the null
model:
$$
 BF(\Mg : \M_0) =    (1 + g)^{(n - 1 - \pg)/2} (1 + g(1 - R^2))^{-(n-1)/2}
$$
\pause
\begin{itemize}
\item Let $g$ be a fixed constant and take $n$ fixed. \pause
\item Let $F = \frac{R_{\g}^2/\pg}{(1 - R_{\g}^2)/(n - 1 - \pg)}$ \pause
\item As $R^2_{\g} \to 1$, $F \to \infty$ LR test would reject $\M_0$
  where $F$ is the usual $F$ statistic for  comparing model $\Mg$ to
  $\M_0$ \pause
\item BF converges to a fixed constant $(1+g)^{-\pg/2}$  (does not go
  to infinity
\end{itemize}

``Information Inconsistency''  see Liang et al JASA 2008


\end{frame}


\begin{frame}
  \frametitle{Mixtures of $g$ priors \& Information consistency}

Need $BF \to \infty$ if $\R^2 \to 1$  $\Leftrightarrow$ $\E_g[(1 +
g)^{-\pg/2}]$ diverges for $\pg < n - 1$ (proof in Liang et al)
\pause
\begin{itemize}
\item Zellner-Siow Cauchy prior \pause
\item hyper-g prior or hyper-g/n (Liang et al JASA 2008) \pause
\item robust prior (Bayarrri et al Annals of Statistics 2012 \pause
\end{itemize}

 All have tails that behave like a Cauchy distribution
\end{frame}
\section{Mortality}
\begin{frame}[fragile]
\frametitle{Mortality \& Pollution}
  \begin{itemize}
  \item Data 60 cities from Statistical Sleuth 12.17 \pause
\item response Mortality \pause
\item 15 predictors;  measures of HC, NOX, SO2 \pause
\item Is pollution associated with mortality after adjusting for other
  socio-economic and meteorological factors? \pause
\item 15 predictor variables implies $2^{15} = 32,768$ possible models
  \end{itemize}
\begin{Schunk}
\begin{Sinput}
> data(ex1217, package="Sleuth3")